% GENERATED FILE. Edit canonical lessons, then run npm run generate:books.
% canonical-source-hash: fnv1a64:48959a428959bac0
% canonical-lessons: ES-C60-mal

\chapter{Mal --- Bien Finally Gets Its Opposite}
\label{ch:mal}
\hlchaptermodality{228}

\begin{chapteropening}
\textbf{By the end of this chapter:} \emph{I can say that something is bad, or going badly.}

I can say that something is bad, or going badly.
\end{chapteropening}

\section[mal]{mal — bien finally gets its opposite}

\label{lesson:ES-C60-mal}



\begin{quote}
\emph{Bien} has been with you since chapter one, and so has \emph{bueno}.

Both have an opposite, and the pair works exactly the same way.
\end{quote}

\begin{cousinweb}
\emph{\textbf{Malo}} is the adjective — \emph{bad}. \emph{\textbf{Mal}} is what is left of it when it goes in front of a masculine noun, and it is also the adverb — \emph{badly}.

\noindent
\begin{tabularx}{\linewidth}{@{}>{\raggedright\arraybackslash}X>{\raggedright\arraybackslash}X@{}}
\toprule
\textbf{} & \textbf{} \\
\midrule
\emph{bien} / \emph{\textbf{mal}} & well / badly \\
\emph{bueno} / \emph{\textbf{malo}} & good / bad \\
\bottomrule
\end{tabularx}

\begin{quote}
\emph{Estoy} \emph{\textbf{mal}}\emph{.} — I'm not well.
\end{quote}

Both come straight from Latin: \emph{bene} and \emph{\textbf{malus}}, which you can still hear in English \emph{malady} and \emph{malice}.
\end{cousinweb}

\begin{culture}
Now look at what has happened twice in three chapters.

\noindent
\begin{tabularx}{\linewidth}{@{}>{\raggedright\arraybackslash}X>{\raggedright\arraybackslash}X>{\raggedright\arraybackslash}X@{}}
\toprule
\textbf{long form} & \textbf{short form} & \textbf{short one appears} \\
\midrule
\emph{mucho} & \emph{\textbf{muy}} & before an adjective or adverb \\
\emph{malo} & \emph{\textbf{mal}} & before a masculine noun, and as the adverb \\
\bottomrule
\end{tabularx}

This is not two exceptions. It is one habit of the language: Spanish \textbf{bites the ending off a word that is leaning on another word}, and keeps it whole when it stands on its own.

Once you have heard it twice you will start noticing it everywhere, and it stops being something to memorise. A short form is a word in a hurry to get to the next one.
\end{culture}

\subsection*{Your turn}
Say these aloud:

\begin{itemize}
\raggedright
  \item \emph{bien}, then \emph{mal}
  \item \emph{mucho} / \emph{muy}, then \emph{malo} / \emph{mal}, and hear the same trick twice
\end{itemize}

\emph{Twice through:}

\begin{tcolorbox}[breakable,colback=teal!4,colframe=teal!35!black,title={Before you move on}]
When does Spanish bite the ending off? (When the word is \textbf{leaning on} another word.)

And which English words still carry Latin \emph{malus}? (\emph{Malady}, \emph{malice}.)
\end{tcolorbox}
